\documentclass[a4paper,12pt]{article}

\usepackage[utf8]{inputenc}
\usepackage[T1]{fontenc}
\usepackage[portuguese]{babel}
\usepackage{amsmath}
\usepackage{booktabs}
\usepackage{graphicx}
\usepackage{geometry}
\geometry{margin=2.5cm}

\title{Aproximação de $\pi$ pela Série de Gregory-Leibniz: \\ Precisão \textit{versus} Custo Computacional}
\author{}
\date{\today}

\begin{document}

\maketitle

\section{Introdução}

Este relatório apresenta a implementação de um programa em linguagem C para calcular
uma aproximação numérica de $\pi$ utilizando a série de Gregory-Leibniz, dada por:

\begin{equation}
    \pi = 4 \sum_{i=0}^{k-1} \frac{(-1)^i}{2i+1}
\end{equation}

O objetivo é analisar como a acurácia da aproximação evolui conforme o número de
iterações $k$ aumenta, relacionando esse ganho de precisão ao custo computacional
(tempo de execução) associado.

\section{Metodologia}

O programa calcula o valor aproximado de $\pi$ para diferentes valores de $k$
(de $10$ a $10^8$), medindo o tempo de execução de cada cálculo com a função
\texttt{clock()} da biblioteca \texttt{time.h}. O erro absoluto foi obtido pela
diferença entre o valor de referência de $\pi$ (com 20 casas decimais) e o valor
calculado pela série.

\section{Resultados}

A Tabela~\ref{tab:resultados} apresenta os resultados obtidos.

\begin{table}[h]
\centering
\caption{Aproximação de $\pi$, erro absoluto e tempo de execução por número de iterações}
\label{tab:resultados}
\begin{tabular}{rrrr}
\toprule
Iterações & $\pi$ calculado & Erro absoluto & Tempo (s) \\
\midrule
10          & 3.0418396189 & $9{,}98 \times 10^{-2}$ & 0.000000 \\
100         & 3.1315929036 & $1{,}00 \times 10^{-2}$ & 0.000000 \\
1.000       & 3.1405926538 & $1{,}00 \times 10^{-3}$ & 0.000000 \\
10.000      & 3.1414926536 & $1{,}00 \times 10^{-4}$ & 0.000000 \\
100.000     & 3.1415826536 & $1{,}00 \times 10^{-5}$ & 0.001000 \\
1.000.000   & 3.1415916536 & $1{,}00 \times 10^{-6}$ & 0.005000 \\
10.000.000  & 3.1415925536 & $1{,}00 \times 10^{-7}$ & 0.051000 \\
100.000.000 & 3.1415926436 & $1{,}00 \times 10^{-8}$ & 0.522000 \\
\bottomrule
\end{tabular}
\end{table}

\section{Análise}

\subsection{Convergência do erro}

Os dados confirmam o comportamento teórico esperado da série de Gregory-Leibniz:
o erro absoluto decresce de forma proporcional a $1/k$. A cada aumento de uma
ordem de grandeza no número de iterações, o erro diminui aproximadamente pela
mesma proporção (uma ordem de grandeza), como pode ser observado diretamente na
tabela: passando de $k=10$ para $k=10^8$, o erro cai de $\sim 10^{-1}$ para
$\sim 10^{-8}$, um ganho de sete ordens de grandeza acompanhando sete ordens de
grandeza no número de iterações.

Essa é uma convergência considerada \emph{lenta} quando comparada a outras
séries para $\pi$ (como as baseadas na fórmula de Machin ou no algoritmo de
Chudnovsky, que convergem exponencialmente). Por isso, a série de
Gregory-Leibniz é didaticamente interessante, mas pouco utilizada na prática
para cálculos de alta precisão.

\subsection{Custo computacional}

O tempo de execução cresce de forma aproximadamente linear com o número de
iterações, como esperado de um algoritmo de complexidade $O(k)$: multiplicar
$k$ por 10 multiplica o tempo de execução por um fator próximo de 10 (por
exemplo, de $10^7$ para $10^8$ iterações, o tempo passa de 0{,}051\,s para
0{,}522\,s).

\subsection{Relação entre precisão e custo}

Combinando as duas observações anteriores, chega-se a um ponto central deste
experimento: \textbf{cada ordem de grandeza adicional de precisão exige uma
ordem de grandeza adicional de processamento}. Diferentemente de algoritmos com
convergência mais rápida, aqui não há "atalho": obter uma casa decimal a mais
de exatidão custa, em média, dez vezes mais tempo de computação. Esse é um
exemplo claro de \emph{retornos decrescentes} do ponto de vista prático: o
esforço computacional cresce muito mais rápido do que o ganho de informação
útil (dígitos corretos) obtido.

\section{Reflexão: aplicações reais}

Esse padrão --- mais processamento gerando ganhos de precisão cada vez mais
caros --- se repete em diversas áreas da computação científica e da
inteligência artificial:

\begin{itemize}
    \item \textbf{Simulações físicas:} em métodos numéricos para equações
    diferenciais (dinâmica de fluidos, clima, mecânica estrutural), refinar a
    malha espacial ou reduzir o passo de tempo melhora a precisão do resultado,
    mas aumenta o custo computacional de forma não linear, frequentemente
    exigindo o uso de supercomputadores e paralelização massiva.

    \item \textbf{Inteligência artificial:} no treinamento de modelos de
    aprendizado de máquina, aumentar o número de dados, épocas ou parâmetros do
    modelo geralmente melhora a acurácia, mas segue uma curva de retornos
    decrescentes semelhante à observada aqui: dobrar o poder computacional não
    dobra o desempenho do modelo. Isso motiva a busca por algoritmos e
    arquiteturas mais eficientes, em vez de depender apenas de mais poder de
    processamento (analogamente ao uso de séries de convergência mais rápida no
    lugar da série de Gregory-Leibniz).

    \item \textbf{Métodos numéricos em geral:} a escolha de um algoritmo com
    melhor taxa de convergência (como Newton-Raphson em vez de bissecção, ou
    Chudnovsky em vez de Leibniz) costuma ser mais eficaz do que simplesmente
    aumentar o número de iterações de um método lento.
\end{itemize}

\section{Conclusão}

O experimento confirma que a série de Gregory-Leibniz converge para $\pi$ de
forma lenta e previsível, com erro proporcional a $1/k$ e tempo de execução
proporcional a $k$. Isso evidencia, de forma simples e mensurável, um princípio
recorrente em aplicações computacionais que demandam alta precisão: ganhos de
exatidão têm custo computacional crescente, e a escolha do algoritmo (não
apenas do volume de processamento) é decisiva para viabilizar resultados
precisos dentro de restrições práticas de tempo e recursos.

\end{document}